% Dydžio ir pastangų vertinimas (atsakingas Marijus Planciunas).
% Lentelių skaičiai – iš galutinio įverčio (2026-09-24). Proza dar nerašyta:
% kiekvienoje vietoje, pažymėtoje [Rengiama], komentaras „TEKSTAS“ nurodo, ką parašyti.

\section{Dydžio ir pastangų vertinimas}\label{sec:dydis}

% TEKSTAS (santraukos pastraipa be poskyrio antraštės, 2–4 sakiniai):
% - dydis 381 UFP (87 funkcijos), AFP ≈ 404; darbo sąnaudos 739,25 žm. d. = 35,2 žm. mėn.;
% - pasiskirstymas: programavimas 309, priedas 231,75, neprograminiai darbai 198,5 žm. d.;
%   su 26 žm. d. rizikų rezervu bazinis planas 765,25 žm. d. (36,4 žm. mėn.);
% - didžiausios eilutės: N2 turinio gamyba (69 žm. d.) ir S1 vykdymo variklis (44,5 žm. d.);
%   išplėtimai (poreikiai 14–20) neįtraukti;
% - čia pirmą kartą apibrėžti FP, žmogaus darbo dienos (žm. d.), žmogaus darbo mėnesiai (žm. mėn.)
%   ir darbo dienos (d. d.); tvarkaraščio skyriuje jų antrą kartą neapibrėžti.
\textit{[Rengiama]}

\subsection{Vertinimo metodas}

Sistemos dydį vertinome funkciniais taškais (FP), o darbo sąnaudas –
kiekvienai sistemos daliai atskirai PERT metodu.

PROBE metodas šiam projektui netinka: jis remiasi komandos istoriniais
duomenimis apie anksčiau sukurtų komponentų dydį ir jų kūrimo trukmę, o tokių
duomenų neturime. FP metodas tinka, nes funkciniai taškai skaičiuojami pagal
tai, ką sistema daro naudotojams, o tai jau žinoma iš poreikių aprašo ir
sistemos vizijos. Programos kodo tam nereikia.

Darbo sąnaudų iš FP tiesiogiai neskaičiavome, nes nežinome, kiek darbo mūsų
komandai reikia vienam funkciniam taškui. Be to, FP neapima vidinio sistemos
sudėtingumo: realiojo laiko vykdymo variklio, grotuvo programos, garso įrangos
integracijos ir infrastruktūros. Todėl kiekvienos dalies sąnaudas įvertinome
ekspertiškai trimis reikšmėmis (optimistine, tikėtiniausia ir pesimistine),
apskaičiavome PERT įvertį ir dalių įverčius sudėjome. FP naudojome dalims
palyginti: didesnei daliai skyrėme daugiau darbo dienų.


\subsection{Sistemos dydis funkciniais taškais}

Sistemos dydis – 381 nekoreguotas funkcinis taškas (UFP). Suskaičiavome
87 funkcijas: 9 vidinius loginius failus (ILF), 47 išorines įvestis (EI),
8 išorines išvestis (EO) ir 23 išorines užklausas (EQ). Išorinių sąsajų
failų (EIF) nėra, nes bazinė sistemos versija su kitomis sistemomis
nesijungia. Funkcijų sudėtingumą nustatėme IFPUG metodo matricomis
(\ref{tab:fp}~lentelė).

Funkciniai taškai skaičiuojami iš naudotojo pusės: kiekviena funkcija,
kurią naudotojas atlieka ar mato, ir kiekviena duomenų grupė, kurią sistema
saugo. Todėl didžiausia dalis yra užsiėmimų programų valdymas (P1, 94~UFP).
Joje treneris kuria, keičia ir šalina programas, jų etapus ir pratimus,
peržiūri programų sąrašą, pratimų katalogą ir versijų istoriją.

Vaizdinio turinio pateikimas (P5) ir garso komunikacija (P6) turi tik 7 ir
8~UFP, nors techniškai yra sudėtingiausios. Naudotojui čia matomos tik kelios
funkcijos, o sinchronizacija realiuoju laiku, grotuvo programa ir radijo
siųstuvo integracija vyksta sistemos viduje, ir FP jų nemato. Šį darbą
vertinome atskirai kaip FP nematomą programinę įrangą (S1--S4,
\ref{tab:pert}~lentelė).


\begin{table}[H]
  \centering
  \caption{Nekoreguoti funkciniai taškai (UFP) pagal sistemos dalis ir funkcijų tipus}
  \label{tab:fp}
  \small
  \begin{tabularx}{\textwidth}{@{}lXrrrrrrr@{}}
    \toprule
    \textbf{Kodas} & \textbf{Sistemos dalis} & \textbf{ILF} & \textbf{EIF} & \textbf{EI} & \textbf{EO} & \textbf{EQ} & \textbf{Funkcijų sk.} & \textbf{UFP} \\
    \midrule
    P0 & Paskyros, sporto erdvės ir įranga & 21 & 0 & 41 & 0 & 23 & 23 & 85 \\
    P1 & Užsiėmimų programų valdymas & 24 & 0 & 45 & 0 & 25 & 22 & 94 \\
    P2 & Užsiėmimų valdymas & 10 & 0 & 40 & 10 & 6 & 13 & 66 \\
    P3 & Dalyvių valdymas & 7 & 0 & 18 & 12 & 18 & 13 & 55 \\
    P4 & Asmeninių programų valdymas & 10 & 0 & 30 & 0 & 10 & 10 & 50 \\
    P5 & Vaizdinio turinio pateikimas & 0 & 0 & 0 & 7 & 0 & 1 & 7 \\
    P6 & Garso komunikacija & 0 & 0 & 3 & 5 & 0 & 2 & 8 \\
    P7 & Duomenų ir istorijos valdymas & 0 & 0 & 0 & 10 & 6 & 3 & 16 \\
    \midrule
    & \textbf{Iš viso, UFP} & \textbf{72} & \textbf{0} & \textbf{177} & \textbf{44} & \textbf{88} & & \textbf{381} \\
    & \textbf{Funkcijų skaičius} & 9 & 0 & 47 & 8 & 23 & \textbf{87} & \\
    \bottomrule
  \end{tabularx}
\end{table}

\subsection{Vertės koregavimo koeficientas}

Nekoreguotą dydį koregavome pagal 14 bendrųjų sistemos charakteristikų
(GSC), kurių kiekviena įvertinta balu nuo 0 iki 5 (\ref{tab:gsc}~lentelė).
Balų suma (TDI) yra 41, todėl vertės koregavimo koeficientas
$\mathrm{VAF} = 0{,}65 + 0{,}01 \times 41 = 1{,}06$, o koreguotas dydis
$\mathrm{AFP} = 381 \times 1{,}06 \approx 404$ funkciniai taškai.

Didžiausią balą (5) gavo tiesioginis duomenų įvedimas: visus duomenis
naudotojai įveda patys per žiniatinklio sąsają, ir jie iškart įrašomi į
sistemą, o paketinio apdorojimo nėra. Dauguma kitų balų ginčytini: jei 12
iš jų būtų vienu balu mažesni ar didesni, AFP svyruotų nuo 358 iki 450.
Darbo sąnaudų įverčiui tai įtakos neturi, nes sąnaudas vertinome kiekvienai
daliai atskirai, o ne iš AFP.


\begin{table}[H]
  \centering
  \caption{Bendrosios sistemos charakteristikos (GSC) ir koreguoti funkciniai taškai}
  \label{tab:gsc}
  \small
  \begin{tabularx}{0.85\textwidth}{@{}rXr@{}}
    \toprule
    \textbf{Nr.} & \textbf{Charakteristika} & \textbf{Įtaka (0--5)} \\
    \midrule
    1 & Duomenų perdavimas & 4 \\
    2 & Paskirstytas duomenų apdorojimas & 4 \\
    3 & Našumas & 3 \\
    4 & Smarkiai apkrauta konfigūracija & 3 \\
    5 & Transakcijų dažnumas & 3 \\
    6 & Tiesioginis duomenų įvedimas & 5 \\
    7 & Efektyvumas galutiniam naudotojui & 3 \\
    8 & Tiesioginis atnaujinimas & 3 \\
    9 & Sudėtingas apdorojimas & 3 \\
    10 & Pakartotinis panaudojimas & 2 \\
    11 & Diegimo lengvumas & 1 \\
    12 & Eksploatavimo lengvumas & 2 \\
    13 & Kelių vietų palaikymas & 3 \\
    14 & Pokyčių palengvinimas & 2 \\
    \midrule
    & \textbf{Bendra įtaka (TDI)} & \textbf{41} \\
    & Vertės koregavimo koeficientas VAF $= 0{,}65 + 0{,}01 \times 41$ & 1,06 \\
    & Koreguoti funkciniai taškai AFP $= 381 \times 1{,}06 = 403{,}86$ & \textbf{$\approx$\,404} \\
    \bottomrule
  \end{tabularx}
\end{table}

\subsection{Darbo sąnaudų įverčiai (PERT)}

Kiekvienai sistemos daliai nurodėme tris darbo sąnaudų įverčius:
optimistinį (O), tikėtiniausią (M) ir pesimistinį (P). Iš jų apskaičiavome
PERT įvertį $E = (O + 4M + P)/6$ ir standartinį nuokrypį $\sigma = (P - O)/6$
(\ref{tab:pert}~lentelė). Tikėtiniausias įvertis turi didžiausią svorį, nes
jis realistiškiausias, o $\sigma$ rodo, kiek neapibrėžtas dalies įvertis.

Pavyzdžiui, dalyvių valdymo (P3) įverčiai yra 18, 25 ir 38~žm.~d., todėl
$E = (18 + 4 \times 25 + 38)/6 = 26$~žm.~d., o $\sigma = (38 - 18)/6
\approx 3{,}3$~žm.~d.

FP matomų sistemos dalių (P0--P7) programavimas sudaro 181~žm.~d., FP
nematomos programinės įrangos (S1--S4) – 128~žm.~d., iš viso 309~žm.~d.

Neprograminius darbus (N1--N5: įrangą, turinio gamybą, bandymus objekte,
mokymus ir dokumentaciją, pilotinį naudojimą ir perdavimą) taip pat
vertinome PERT metodu; jie sudaro 198,5~žm.~d. Procentinio priedo prie jų
nepridėjome, nes tai nėra programų realizacija, o dalis šių darbų
(integraciniai bandymai objekte ir priėmimo bandymas) pati yra testavimas.


\begin{table}[H]
  \centering
  \caption{Darbo sąnaudų įverčiai PERT metodu, žm.~d.}
  \label{tab:pert}
  \small
  \setlength{\tabcolsep}{5pt}
  \begin{tabularx}{\textwidth}{@{}l>{\raggedright\arraybackslash}Xrrrrrr@{}}
    \toprule
    \textbf{Kodas} & \textbf{Eilutė} & \textbf{O} & \textbf{M} & \textbf{P} & \textbf{E} & {\boldmath$\sigma$} & \textbf{\makecell[r]{žm.~d. /\\UFP}} \\
    \midrule
    P0 & Paskyros, sporto erdvės ir įranga & 23 & 33 & 50 & 34,0 & 4,5 & 0,40 \\
    P1 & Užsiėmimų programų valdymas & 27 & 38 & 58 & 39,5 & 5,2 & 0,42 \\
    P2 & Užsiėmimų valdymas & 22 & 32 & 50 & 33,5 & 4,7 & 0,51 \\
    P3 & Dalyvių valdymas & 18 & 25 & 38 & 26,0 & 3,3 & 0,47 \\
    P4 & Asmeninių programų valdymas & 15 & 22 & 34 & 23,0 & 3,2 & 0,46 \\
    P5 & Vaizdinio turinio pateikimas & 4 & 8 & 14 & 8,5 & 1,7 & 1,21 \\
    P6 & Garso komunikacija & 4 & 7 & 12 & 7,5 & 1,3 & 0,94 \\
    P7 & Duomenų ir istorijos valdymas & 5 & 9 & 14 & 9,0 & 1,5 & 0,56 \\
    \midrule
    & \textbf{A. FP matomos sistemos dalys} & & & & \textbf{181,0} & & \\
    \midrule
    S1 & Realiojo laiko užsiėmimo vykdymo variklis ir sinchronizacija & 28 & 42 & 72 & 44,5 & 7,3 & -- \\
    S2 & Vietinio grotuvo programa & 25 & 37 & 62 & 39,0 & 6,2 & -- \\
    S3 & Garso įrangos integracija & 10 & 18 & 35 & 19,5 & 4,2 & -- \\
    S4 & Infrastruktūra ir saugumas & 15 & 24 & 38 & 25,0 & 3,8 & -- \\
    \midrule
    & \textbf{B. FP nematoma programinė įranga} & & & & \textbf{128,0} & & \\
    & \textbf{Programavimas iš viso (A + B)} & & & & \textbf{309,0} & & \\
    \midrule
    N1 & Techninės įrangos parinkimas, užsakymas ir montavimas & 17 & 26 & 42 & 27,0 & 4,2 & -- \\
    N2 & Pratimų vaizdo ir garso turinio gamyba & 45 & 66 & 105 & 69,0 & 10,0 & -- \\
    N3 & Integraciniai bandymai objekte & 22 & 35 & 60 & 37,0 & 6,3 & -- \\
    N4 & Trenerių ir administratorių mokymai, naudotojo dokumentacija & 16 & 25 & 40 & 26,0 & 4,0 & -- \\
    N5 & Pilotinis naudojimas, priėmimo bandymas, perdavimas ir garantinis palaikymas & 24 & 38 & 62 & 39,5 & 6,3 & -- \\
    \midrule
    & \textbf{C. Neprograminiai darbai} & & & & \textbf{198,5} & & \\
    \bottomrule
  \end{tabularx}
\end{table}

\subsection{Reikalavimų, projektavimo, testavimo ir valdymo priedas}

Programavimo įvertis (309~žm.~d.) neapima kitų projekto veiklų. Reikalavimų
analizės, projektavimo, testavimo ir projekto valdymo sąnaudas apskaičiavome
procentiniu metodu, t.~y. kaip dalį nuo programavimo sąnaudų: reikalavimai
15~\%, projektavimas 20~\%, testavimas 30~\%, projekto valdymas 10~\%. Iš viso
tai sudaro 75~\%, arba 231,75~žm.~d. (\ref{tab:priedas}~lentelė).

Priedas nėra rezervas. Tai suplanuotas darbas, kuris tvarkaraštyje turi
atskiras veiklas: reikalavimų specifikaciją, architektūros projektavimą,
kiekvienos iteracijos testavimą ir projekto valdymą. Nenumatytiems atvejams
skirtas atskiras rizikų rezervas (\ref{tab:rezervas}~lentelė).


\begin{table}[H]
  \centering
  \caption{Priedas prie programavimo sąnaudų (309~žm.~d.) pagal paskaitos procentus}
  \label{tab:priedas}
  \small
  \begin{tabular}{@{}lrr@{}}
    \toprule
    \textbf{Veikla} & \textbf{Dalis, \%} & \textbf{Sąnaudos, žm.~d.} \\
    \midrule
    Reikalavimai & 15 & 46,35 \\
    Projektavimas & 20 & 61,80 \\
    Testavimas & 30 & 92,70 \\
    Projekto valdymas & 10 & 30,90 \\
    \midrule
    \textbf{Iš viso} & \textbf{75} & \textbf{231,75} \\
    \bottomrule
  \end{tabular}
\end{table}

\subsection{Bendros sąnaudos ir rezervai}

% TEKSTAS:
% - suvestinė (\ref{tab:sanaudos}~lentelė): 739,25 žm. d. + 26 rizikų rezervas = bazinis planas
%   765,25 žm. d.; + 37 valdymo rezervas = 802,25 žm. d.;
% - čia pirmą kartą pavartojama EMV, todėl apibrėžti vieną kartą: „tikėtinoji vertė (angl. expected
%   monetary value, EMV)“ – rizikos tikimybė, padauginta iš jos papildomų sąnaudų; toliau tik „EMV“;
% - rizikų rezervas = Σ EMV (25,85 ≈ 26 žm. d., \ref{tab:rezervas}~lentelė); kodėl ne 1,28σ (≈ 38 žm. d.):
%   σ rodo įverčių sklaidą, o registras – atskirus įvykius; susieti su rizikų valdymo skyriaus registru;
% - rizikų rezervą valdo projekto vadovas, valdymo rezervą (5~\%) laiko užsakovas; apimties pokyčiai
%   tvirtinami per pakeitimų valdymą, ne iš rezervų;
% - rezervai skirti sąnaudoms, ne kalendoriui (laiko rezervai – \ref{sec:tvarkarastis} skyriuje).
\textit{[Rengiama]}

\begin{table}[H]
  \centering
  \caption{Bendros darbo sąnaudos ir rezervai (1~žm.~mėn. = 21~d.~d.)}
  \label{tab:sanaudos}
  \small
  \begin{tabular}{@{}lrr@{}}
    \toprule
    \textbf{Sudedamoji dalis} & \textbf{Žm.~d.} & \textbf{Žm.~mėn.} \\
    \midrule
    Programavimas (A + B) & 309,00 & 14,71 \\
    Priedas (75~\%) & 231,75 & 11,04 \\
    Neprograminiai darbai (C) & 198,50 & 9,45 \\
    \midrule
    \textbf{Iš viso be rezervų} & \textbf{739,25} & \textbf{35,20} \\
    Rizikų rezervas ($\Sigma$ EMV) & 26,00 & 1,24 \\
    \textbf{Bazinis planas} & \textbf{765,25} & \textbf{36,44} \\
    Valdymo rezervas (5~\%, laiko užsakovas) & 37,00 & 1,76 \\
    \midrule
    \textbf{Iš viso su rezervais} & \textbf{802,25} & \textbf{38,20} \\
    \bottomrule
  \end{tabular}
\end{table}

\begin{table}[H]
  \centering
  \caption{Rizikų rezervo apskaičiavimas (EMV = tikimybė $\times$ papildomos sąnaudos)}
  \label{tab:rezervas}
  \small
  \begin{tabularx}{\textwidth}{@{}rl>{\raggedright\arraybackslash}Xrrr>{\raggedright\arraybackslash}p{2.1cm}@{}}
    \toprule
    \textbf{Nr.} & \textbf{ID} & \textbf{Rizika} & \textbf{\makecell[r]{Tikimybė,\\\%}} & \textbf{\makecell[r]{Papildomos\\sąnaudos,\\žm.~d.}} & \textbf{\makecell[r]{EMV,\\žm.~d.}} & \textbf{\makecell[l]{Paveiktos\\eilutės}} \\
    \midrule
    1 & TR1 & Drėgmė ir kondensatas pažeidžia stacionarią įrangą baseine & 35 & 8 & 2,80 & N1, N3, N5 \\
    2 & TR2 & Nepakankamas belaidis ryšys baseino salėje & 40 & 6 & 2,40 & N1, N3, S1, S2 \\
    3 & FR1 & Vėlai aptiktas užsiėmimo būsenų mašinos defektas & 25 & 8 & 2,00 & S1, S2, P2, N3, N5 \\
    4 & IR1 & Grotuvų, ekranų ar sporto salės kolonėlių ribojimai & 35 & 8 & 2,80 & S2, S3, N1, N3 \\
    5 & IR2 & Neaiški radijo siųstuvo ir trenerio mikrofono sujungimo grandinė & 30 & 8 & 2,40 & S2, S3, P6, N1, N3 \\
    6 & OR1 & Treneriams sunku perimti užsiėmimo valdymą & 25 & 5 & 1,25 & P2, N4, N5 \\
    7 & SR1 & Griežtesni asmens duomenų reikalavimai arba prieigos teisių spraga & 30 & 7 & 2,10 & P0, P3, P7, S4 \\
    8 & ER1 & Pratimų vaizdas prastai matomas iš tolimų takelių & 40 & 10 & 4,00 & N2, P5, S2, N1, N3 \\
    9 & ER2 & Ausinės slysta nuo kepurėlių arba yra nepatogios & 20 & 5 & 1,00 & N1, N3, S3 \\
    10 & OR2 & Įranga vėluoja arba atkeliauja kitas modelis ar mikroprograma & 30 & 7 & 2,10 & S2, S3, N1, N3 \\
    11 & OR3 & Dalies pratimų perfilmavimas ir garso instrukcijų įrašymas iš naujo & 30 & 10 & 3,00 & N2, N5 \\
    \midrule
    & & \textbf{Iš viso} & & \textbf{82} & \textbf{25,85\,$\approx$\,26} & \\
    \bottomrule
  \end{tabularx}
\end{table}

% TEKSTAS (po \ref{tab:rezervas}~lentele, 1–2 sakiniai):
% - ID sutampa su rizikų valdymo skyriaus registru (Godos); IR2, OR2 ir OR3 – naujos rizikos;
%   debesijos paslaugų teikėjo sutrikimas (TR3) į rezervą neįtrauktas, nes rizika perduota
%   paslaugų teikėjui;
% - registre nurodyta pradinė tikimybė (1–5 balai: 1 = 0–19~\%, 2 = 20–39~\%, 3 = 40–59~\%,
%   4 = 60–79~\%, 5 = 80–100~\%), o lentelėje – likutinė tikimybė po tvarkaraštyje suplanuotų
%   mažinimo veiklų: TR1 4 balai → 35~\% po bandomojo komplekto bandymo baseine,
%   IR1 3 balai → 35~\% po prototipo bandymo baseine, ER1 4 balai → 40~\% po bandomojo
%   klipo matomumo patikros baseine prieš montuojant visus pratimus;
% - „Kitų registro rizikų (TR2, FR1, OR1, SR1, ER2) procentai atitinka Godos balus; naujų IR2, OR2
%   ir OR3 – po 30~\% (2 balai).“;
% - prozoje veiklų ID (A10, A21, A74) nerašyti: rašyti veiklos pavadinimą, kaip nurodyta šiame
%   komentare, arba nuorodą į \ref{sec:tvarkarastis} skyrių.
\textit{[Rengiama]}

\subsection{Prielaidos}

% TEKSTAS (sąrašas mažąja raide, punktai skiriami kabliataškiais):
% - vertinama bazinė versija (poreikiai 1–13), vienas sporto centras: baseinas ir sporto salė;
%   treneriai vaizdo įrašų nekelia, 4K nėra, turinys filmuojamas iš anksto (Tado sprendimas);
% - pradinė biblioteka – apie 40 pratimų ir 10 programų; architektūra: NestJS, PostgreSQL, Redis,
%   React, vietiniai DSDX5 grotuvai;
% - programuotojai dirba visu etatu, vidutinio lygio; atostogos vertinamos tvarkaraštyje, ne sąnaudose.
\textit{[Rengiama]}
